\documentclass[12pt]{article}
\usepackage[margin=1in]{geometry}
\usepackage{titlesec}
\usepackage{hyperref}
\usepackage{enumitem}
\usepackage{parskip}
\usepackage{booktabs}
\usepackage{array}
\usepackage{fancyhdr}
\usepackage{longtable}

\pagestyle{fancy}
\fancyhf{}
\rhead{CS 3338 --- Group 11}
\lhead{Scientific Forgery Image Detection}
\cfoot{\thepage}

\titleformat{\section}{\large\bfseries}{\thesection.}{0.5em}{}[\titlerule]
\titleformat{\subsection}{\normalsize\bfseries}{\thesubsection}{0.5em}{}

\hypersetup{
    colorlinks=true,
    linkcolor=blue,
    urlcolor=blue
}

\begin{document}

% Title Page
\begin{center}
    \vspace*{2cm}
    {\LARGE \textbf{Scientific Forgery Image Detection}} \\[1em]
    {\large By: Janelle Rivera, Eric Marroquin, Alexis Flores, Alex Lam} \\[0.5em]
    {\large May 8, 2026} \\[0.5em]
    {\large Software Engineering Tools, CS 3338}
\end{center}

\vspace{2cm}

\tableofcontents
\newpage

% Section 1: Introduction
\section{Introduction}

\subsection{Purpose}

Scientific images play an important role in academic research and published studies. However, some images may be manipulated through techniques such as copy-move forgery, where parts of an image are duplicated and placed elsewhere to falsify results.

The purpose of this system is to detect and analyze copy-move forgery in biomedical and scientific images. Users will be able to upload images through a web interface, and the system will process them to identify potentially duplicated regions. This helps support scientific integrity by assisting researchers and reviewers in identifying manipulated figures.

\subsection{Intended Audience}

This document is intended for:
\begin{itemize}[noitemsep]
    \item Software developers working on the system
    \item Course instructors evaluating the project
    \item Students involved in system design
    \item Future maintainers of the software
\end{itemize}

\subsection{Overview}

The system is a web-based application consisting of a frontend interface and a backend server. Users upload images through the frontend, and the backend processes these images for future detection logic development.

At this stage, the system focuses on:
\begin{itemize}[noitemsep]
    \item Project setup
    \item Frontend image upload structure
    \item Backend API structure
    \item Fraudulent region detection workflow using OpenCV
    \item Basic results display with highlighted suspicious regions
\end{itemize}

% Section 2: System Architecture
\section{System Architecture}

\subsection{Overview}

The system follows a simple client-server architecture.

\subsection{Frontend}

The frontend is responsible for user interaction.

\textbf{Responsibilities:}
\begin{itemize}[noitemsep]
    \item Provide image upload interface
    \item Send image data to backend
    \item Display backend responses
    \item Show suspicious regions
\end{itemize}

\textbf{Technologies:}
\begin{itemize}[noitemsep]
    \item React.js
    \item HTML/CSS
\end{itemize}

\subsection{Backend}

The backend handles incoming requests from the frontend.

\textbf{Responsibilities:}
\begin{itemize}[noitemsep]
    \item Receive uploaded images
    \item Validate image file format
    \item Run fraudulent region detection
    \item Return placeholder or future detection results
\end{itemize}

\textbf{Technologies:}
\begin{itemize}[noitemsep]
    \item Python
    \item FastAPI
    \item Uvicorn
    \item OpenCV
    \item NumPy
    \item Pillow
\end{itemize}

\subsection{Detection Module}

The detection module analyzes uploaded images for copy-move forgery.

\textbf{Responsibilities:}
\begin{itemize}[noitemsep]
    \item Detect keypoints in images using ORB
    \item Match keypoints to identify duplicated regions
    \item Highlight suspicious regions on the output image
    \item Return detection results to backend
\end{itemize}

\subsection{Workflow Design}

\textbf{Updated Workflow:}
\begin{enumerate}[noitemsep]
    \item User uploads biomedical image through frontend
    \item Frontend sends image to backend via POST request
    \item Backend receives and validates the image
    \item Detection module analyzes the image for duplicated regions
    \item Suspicious regions are identified and marked
    \item Detection results returned to frontend
    \item User views highlighted suspicious regions
\end{enumerate}

% Section 3: User Interface Design
\section{User Interface Design}

\subsection{Homepage}

The homepage introduces the application and provides navigation to upload functionality.

\textbf{Elements:}
\begin{itemize}[noitemsep]
    \item Project title
    \item Short description
    \item Navigation to upload page
\end{itemize}

\subsection{Image Upload Interface}

This page allows users to submit images for analysis.

\textbf{Elements:}
\begin{itemize}[noitemsep]
    \item Upload button
    \item File selection option
    \item Submit button
    \item Image preview
\end{itemize}

\subsection{Results Display}

\begin{itemize}[noitemsep]
    \item Confirmation of upload
    \item Detection status message
    \item Marked image showing highlighted suspicious regions
    \item Match count indicator
\end{itemize}

% Section 4: Glossary
\section{Glossary}

\begin{longtable}{>{\bfseries}p{0.28\textwidth} p{0.67\textwidth}}
\toprule
Term & Definition \\
\midrule
\endhead
API & Application Programming Interface used for communication between frontend and backend \\[0.5em]
Backend & Server-side system that processes requests \\[0.5em]
Frontend & User interface of the application \\[0.5em]
Copy-Move Forgery & Image manipulation where parts of an image are duplicated within the same image \\[0.5em]
React.js & JavaScript library used for building user interfaces \\[0.5em]
FastAPI & Python framework used for building APIs \\[0.5em]
Uvicorn & Server used to run FastAPI applications \\[0.5em]
OpenCV & Open-source computer vision library used for image analysis \\[0.5em]
ORB & Oriented FAST and Rotated BRIEF --- keypoint detection algorithm \\[0.5em]
Keypoint & A distinctive location in an image used for matching \\
\bottomrule
\end{longtable}

% Section 5: Versions Table
\section{Versions Table}

\begin{table}[h]
\centering
\begin{tabular}{>{\centering}p{0.1\textwidth} >{\centering}p{0.15\textwidth} p{0.38\textwidth} p{0.25\textwidth}}
\toprule
\textbf{Version} & \textbf{Date} & \textbf{Description} & \textbf{Authors} \\
\midrule
1.0 & May 8, 2026 & Initial Software Design Document created for Snapshot 1 & Janelle Rivera, Eric Marroquin, Alexis Flores \\[0.5em]
1.1 & May 10, 2026 & Updated frontend implementation progress and UI design & Janelle Rivera, Eric Marroquin, Alexis Flores, Alex Lam \\[0.5em]
1.2 & May 12, 2026 & Added fraudulent region detection workflow and frontend/backend integration updates & Janelle Rivera, Eric Marroquin, Alexis Flores, Alex Lam \\
\bottomrule
\end{tabular}
\end{table}

% Section 6: References
\section{References}

\begin{itemize}[noitemsep]
    \item FastAPI Documentation: \url{https://fastapi.tiangolo.com/}
    \item React Documentation: \url{https://react.dev/}
    \item OpenCV Documentation (future reference): \url{https://opencv.org/}
\end{itemize}

\end{document}